% ============================================================
% 《从一个窗口到一个引擎》—— C++ 游戏引擎与可视化编辑器构建实录
% 主文件：ElegantBook 模板 + XeLaTeX 编译
% ============================================================
\documentclass[%
  lang=cn,          % 中文支持（ctex）
  newtx,            % Times 风格字体
  12pt,             % 教材正文偏大字号
  oneside,          % 单面排版（屏幕阅读）
  scheme=warmgreen, % 主题色
]{elegantbook}

% ---------- 项目信息 ----------
\title{从一个窗口到一个引擎}
\subtitle{C++ 游戏引擎与可视化编辑器构建实录}
\author{LAJENUNESSE}
\date{\today}

% ---------- 全局代码样式 ----------
% 引擎代码为 C++20，中文注释；listings 对 UTF-8 中文注释的支持有限，
% 故代码清单中的注释以英文呈现，讲解文字承载中文说明。
\lstset{
  language=C++,
  basicstyle=\small\ttfamily,
  keywordstyle=\bfseries,
  commentstyle=\color{gray!70}\itshape,
  stringstyle=\color{teal},
  numbers=left,
  numberstyle=\tiny\color{gray},
  numbersep=8pt,
  frame=single,
  breaklines=true,
  tabsize=4,
  showstringspaces=false,
  columns=flexible,
  keepspaces=true,
  escapeinside={(*@}{@*)},
  morekeywords={uint32_t,uint64_t,Scope,Ref,override,consteval,constinit,concept,requires}
}

% 注：ElegantBook 已内置加载 natbib，此处不再重复引入

\begin{document}

\frontmatter
\maketitle

\begin{introduction}
  \item 本书以一个真实完成的毕业设计项目为载体——约 4.6 万行 C++20 代码的
        3D 游戏引擎与可视化编辑器——按照它实际被构建出来的顺序逐章展开，
        讲述每项技术\emph{为什么被引入}、\emph{有哪些备选方案}、
        \emph{最终如何实现}，以及\emph{作者踩过哪些坑}。
  \item 全书不是教程式的理想化重演：所有代码锚点均指向仓库当前的真实文件，
        所有设计决策均有 git 提交历史或项目文档背书。
  \item 关键词：游戏引擎；C++20；ECS；OpenGL；Vulkan；反射系统；可视化编辑器
\end{introduction}

\tableofcontents

\mainmatter

% ---- 第 I 部分：地基 ----
\chapter{绪论：这本书讲什么、怎么读}
\label{ch:intro}
% 第 1 章 绪论（占位——M3 阶段扩写）
本书记录一个真实项目的构建历程。项目始于 2026 年 2 月 21 日，首个提交
包含窗口、主循环与 ImGui 视口；至 2026 年 8 月，仓库已积累 490 次提交，
涵盖渲染（OpenGL/Vulkan 双后端）、ECS 场景、反射驱动编辑器、物理、SPH
流体、CUDA 加速等子系统。

本章将交代全书的组织方式：按构建历程而非架构分层展开；每章固定走
``问题引入 → 通用原理 → 设计决策 → 实现走读 → 踩坑记录''五段结构；
所有代码锚点指向仓库当前真实文件。

% TODO(M3): 扩写项目全景图（TikZ 架构总览）、阅读路线建议、环境搭建


\chapter{引擎骨架：Application 与主循环}
\label{ch:core-loop}
% ============================================================
% 第 2 章 引擎骨架：Application 与主循环（样章）
% 代码锚点：Engine/src/Core/Application.h / Application.cpp
%           Engine/src/Core/LayerStack.h
% 时间线：2026-02-21 首个提交即包含 Application 骨架
% ============================================================
\section{问题引入：一个窗口背后需要什么}

一切从最小的问题开始：在屏幕上打开一个窗口，让它响应鼠标拖动、键盘输入，
并且不卡死操作系统。裸写 Win32 API 需要 200 余行样板代码，且与平台强绑定。
但"跨平台窗口"只是表象——真正的需求是：\textbf{引擎的其余所有子系统
（渲染、物理、音频）都必须挂在一个统一的节拍器上运转}。这个节拍器就是主循环。

本项目选择 GLFW 作为窗口库，但它只解决"窗口 + 输入事件"，不解决：
帧与帧之间该按什么顺序更新？谁拥有这些子系统？窗口关闭时如何优雅退出？
这些问题由 \texttt{Application} 类回答。

\section{通用原理：游戏循环的三种形态}

教科书上的游戏循环通常有三种实现策略：

\begin{enumerate}
  \item \textbf{自由奔跑}：能跑多快跑多快，用帧间隔作为时间步。简单，
        但不同机器上模拟速度不一致，且 GPU 空转烧电。
  \item \textbf{固定时间步}：以固定 $\Delta t$ 累积更新，渲染插值。
        物理仿真的黄金标准，但实现复杂度高。
  \item \textbf{限帧混合睡眠}：目标帧率下，先粗粒度 sleep 让出 CPU，
        最后几毫秒自旋等待以提高唤醒精度。
\end{enumerate}

本项目的选择是第 3 种——这个决策值得玩味：它不是学术最优解，而是工程
权衡解。物理系统（Bullet3）内部自带固定步长累积器，所以主循环无需再背
这份责任；而编辑器场景下"144 FPS 上限 + 低 CPU 占用"比"完美确定性"
重要得多。

\section{本项目的设计与决策}

\subsection{单例还是依赖注入？}

\texttt{Application} 采用了经典的 Meyer's 单例模式：

\begin{lstlisting}
class Application
{
public:
    static Application& Get() { return *s_Instance; }
    // ...
private:
    static Application* s_Instance;
};

// To be defined by the client application
Scope<Application> CreateApplication();
\end{lstlisting}

这是从 Hazel 引擎教学项目继承的模式。它的争议点在于全局状态不利于测试，
但对引擎这种"天然进程唯一"的对象，单例换来的访问便利性是划算的。
值得注意的是 \texttt{CreateApplication()} 自由函数：\emph{引擎不知道客户端是谁}，
由客户端 exe（Editor 或 Sandbox）提供入口——这是控制反转的最小实践。

\subsection{层栈：编辑器的组织单元}

\texttt{LayerStack}（\texttt{Core/LayerStack.h}）把每一帧的工作切分为若干
\texttt{Layer}，每层有 \texttt{OnAttach/OnUpdate/OnImGuiRender/OnEvent/OnDetach}
五个生命周期钩子。Overlay（如 ImGui 层）永远最后更新、最先收到事件。
编辑器本体就是一个被推入栈中的 \texttt{EditorLayer}。

\subsection{帧边界的显式化：为 Vulkan 埋下的伏笔}

最初的 \texttt{Run()} 循环非常直白：更新层 → 渲染 ImGui → 交换缓冲。
但半年后的 Vulkan 迁移打破了这个宁静——Vulkan 的命令录制必须发生在
"acquire swapchain image 之后、submit 之前"，而 OpenGL 没有这个约束。
最终的主循环长这样：

\begin{lstlisting}
void Application::Run()
{
    while (m_Running)
    {
        // ... timestep 计算, PerformanceMonitor::BeginFrame ...

        GraphicsContext* gfx = m_Window->GetGraphicsContext();
        if (gfx)
            gfx->BeginRenderFrame();   // Vulkan: acquire + Begin cmd

        if (!m_Minimized)
        {
            for (const auto& layer : m_LayerStack)
                layer->OnUpdate(timestep);   // 录制 dispatch 到主帧 cmd

            m_ImGuiLayer->Begin();
            for (const auto& layer : m_LayerStack)
                layer->OnImGuiRender();
            m_ImGuiLayer->End();
        }

        if (gfx)
            gfx->EndRenderFrame();     // Vulkan: submit + present

        m_Window->OnUpdate();
        PerformanceMonitor::Get().EndFrame();
        // ... 混合睡眠限帧 ...
    }
}
\end{lstlisting}

关键设计：\texttt{GraphicsContext} 基类的 \texttt{BeginRenderFrame/
EndRenderFrame} 是\textbf{默认 no-op 的虚函数}\footnote{对应迁移文档
docs/vulkan-migration/SPEC.md 决策 D-16。}。OpenGL 路径完全不受影响，
Vulkan 路径 override 它们来驱动帧边界。这是"抽象层吸收后端差异"原则
在整个项目中最典型的一次应用——代价只是两个空虚函数。

\section{实现走读：那些容易忽略的细节}

\subsection{关闭拦截器：一行代码背后的产品思维}

\begin{lstlisting}
using CloseInterceptFn = std::function<bool()>;
void SetCloseInterceptor(CloseInterceptFn fn) { /*...*/ }

bool Application::OnWindowClose(WindowCloseEvent& e)
{
    if (m_CloseInterceptor && !m_CloseInterceptor())
    {
        glfwSetWindowShouldClose(/*...*/, GLFW_FALSE);
        return true;   // 取消关闭
    }
    m_Running = false;
    return true;
}
\end{lstlisting}

用户点击编辑器右上角 × 时，如果有未保存的场景，应该弹窗询问而非直接退出。
这个"产品级细节"通过回调注入实现：引擎核心不知道"未保存"是什么，
编辑器层注册拦截器自行判断。\texttt{OnWindowClose} 里直接调用 GLFW 而非
走抽象层的细节也值得一提——因为取消关闭本质上是窗口操作，强行抽象反而
增加无收益的间接层。

\subsection{析构顺序：注释里的事故现场}

\begin{lstlisting}
Application::~Application()
{
    // 先 detach 所有层——OnDetach 可能需要 Application::Get()
    //（如清除 CloseInterceptor）
    m_LayerStack.Shutdown();

    if (m_Initialized) { /* PerformanceMonitor / AssetManager */ }
    Renderer::Shutdown();
    s_Instance = nullptr;
}
\end{lstlisting}

这行注释是真实的踩坑记录：曾因析构顺序错误导致 \texttt{Application::Get()}
返回悬垂指针崩溃（git 提交 \texttt{2c5ad95}``修复关闭窗口时空指针崩溃''）。
\textbf{教训：带单例的系统里，销毁顺序必须保证"使用单例者先于单例本身死亡"}。

\subsection{限帧的实现：为什么最后 2ms 要自旋}

\begin{lstlisting}
// 混合睡眠策略：先粗粒度 sleep 让出 CPU，
// 最后 2ms 用 _mm_pause 自旋等待以提高精度
if (m_FrameRateLimitEnabled && !m_Window->IsVSync() && ...)
{
    double targetTime = m_LastFrameTime + 1.0 / m_TargetFrameRate;
    while (glfwGetTime() < targetTime - 0.002)
        std::this_thread::sleep_for(std::chrono::milliseconds(1));
    while (glfwGetTime() < targetTime)
    {
        _mm_pause();   // x86；ARM64 对应 __yield()
    }
}
\end{lstlisting}

纯 \texttt{sleep(1ms)} 的唤醒精度受 Windows 定时器分辨率制约（默认
15.6ms），实测会造成帧率抖动；纯自旋则把一个核烧满。混合方案在两者间
取得平衡——注意构造函数里的 \texttt{timeBeginPeriod(1)}：没有它，
Windows 下 \texttt{sleep(1)} 实际睡约 15ms。这个细节在 git 历史中对应
提交 \texttt{9eb9a79}``spin-wait 帧率限制器 CPU 空转 56\% $\to$ ~0\%''。

\section{踩坑记录}

\begin{description}
  \item[Windows.h 宏污染] \texttt{EntryPoint.h} 必须
        \texttt{\#define NOMINMAX}，否则 \texttt{windows.h} 的
        \texttt{min/max} 宏会击穿所有 \texttt{std::min/max} 调用
        （提交 \texttt{8d67964}）。Windows 平台开发的第一课，
        但几乎每个人都要亲自中招一次。
  \item[float 计时的精度陷阱] 帧计时成员最初是 \texttt{float}，
        引擎运行数十分钟后秒级时间戳的 float 精度不足导致 timestep
        异常，后改为 \texttt{double}（提交 \texttt{80fccb2}）。
        "时间相关字段一律 double"由此成为项目约定。
\end{description}

\section{小结与延伸思考}

本章的 \texttt{Application} 只有约 200 行，却确立了三个贯穿全书的原则：
\begin{itemize}
  \item \textbf{引擎不认识客户端}——控制反转让 Editor/Sandbox 复用同一内核；
  \item \textbf{层栈是一切的容器}——后续的功能扩展几乎都以新 Layer 或
        Layer 内控制器形式出现；
  \item \textbf{抽象层的虚函数可以很便宜}——no-op 默认实现是吸收后端差异
        的低成本手段，Vulkan 帧边界改造证明了这一点。
\end{itemize}

延伸思考：如果重做一次，是否会选固定时间步 + 渲染插值？答案可能是否定
的——编辑器的交互流畅度优先级高于模拟确定性，且 Bullet 内部已有自己的
固定步长机制。"在正确的层级解决问题"比"用最先进的方案"更重要。


% ---- 第 II 部分：看见画面 ----
\chapter{渲染抽象层：在 OpenGL 之上筑墙}
\label{ch:render-abstraction}
% 第 3 章 渲染抽象层（占位）
锚点：\texttt{RendererAPI.h}（纯虚接口）、\texttt{RenderCommand.h}
（静态转发门面 + 性能统计内联）、\texttt{Buffer.h/Shader.h/Texture.h/
Framebuffer.h/VertexArray.h} 工厂族、\texttt{Platform/OpenGL/} 22 文件实现。

叙事线：项目早期直接调 GL（Phase R1 前后对比，提交 \texttt{5043e83}）；
2026-05 Vulkan 迁移 Phase 1``抽象泄漏归零''——Engine/src 下零
glad include（终态 \texttt{27398fc}），SSAO 噪声/blend 状态/异步回读/
能力查询/IBL 全部下沉 Platform 层。

关键代码：DrawIndexed 内嵌 stats 统计；BindTextureView(void*) 避开
vulkan.h 泄漏的设计；Material 的 Set-缓存 + Bind-刷新 dirty flag 模式。

% TODO(M3): RenderPacket/RenderQueue 按 Shader 排序减状态切换；HDR
% Framebuffer 附件布局；多 Pass 管线总览图（TikZ）。


\chapter{场景与 ECS：EnTT 与 façade 架构}
\label{ch:scene-ecs}
% ============================================================
% 第 4 章 场景与 ECS：EnTT 与 façade 架构
% 锚点：Engine/src/Scene/Scene.h, Components.h, Entity.h,
%       SceneHierarchyService.cpp, WorldTransformService.cpp,
%       Runtime/SceneRuntimeCoordinator, Systems/
% 时间线：8150251(2026-02-21) → 4dab189 → a6676de → 83dcf39
% ============================================================
\section{问题引入：对象之间说不清的关系}

有了渲染能力后，下一个问题是：场景里的东西怎么组织？一个``点光源挂在
移动的汽车上''的需求，就同时牵出了实体表示、组件组合、父子变换传播、
按类型遍历四件事。用面向对象的实体继承树（GameObject $\to$
LightObject/PhysicsObject……）很快会陷入菱形继承泥潭——这正是 ECS
（Entity-Component-System）要解决的问题。

本项目在首个提交日（2026-02-21）就选定了 EnTT 作为 ECS 底座
（提交 \texttt{8150251} ``Phase 2 Scene/ECS 场景系统 + 编辑器''），
此后经历三轮架构演进：运行时资源剥离（\texttt{4dab189}）、façade 收束
（\texttt{a6676de}）、组件元数据统一（\texttt{83dcf39}）。本章沿这条
演化线展开。

\section{通用原理：ECS 与 sparse-set 存储}

ECS 的核心三分法：\emph{实体}是 ID，\emph{组件}是纯数据 POD，
\emph{系统}是操作特定组件组合的逻辑。数据导向设计让缓存命中率成为
一等公民。

EnTT 的存储方案是 sparse-set：每个组件类型维护一个``稀疏数组（实体 ID
$\to$ 稠密下标）+ 稠密数组（紧凑的组件数据）''对。迭代某类组件时只走稠密
数组，内存连续、缓存友好；查询``拥有 A 且 B 的实体''则取两个集合的交集。
代价是非连续访问（如随机查询单个实体）需要一次稀疏跳转。理解这一点，
才能理解为什么本项目的渲染系统全部写成
\texttt{registry.view<TransformComponent, MeshRendererComponent>()}
式的批量遍历。

\section{本项目的设计与决策}

\subsection{组件即 POD}

\begin{lstlisting}
// Components.h —— 组件是纯数据 + 默认/拷贝构造
struct TransformComponent
{
    glm::vec3 Translation = {0.0f, 0.0f, 0.0f};
    glm::vec3 Rotation    = {0.0f, 0.0f, 0.0f};
    glm::vec3 Scale       = {1.0f, 1.0f, 1.0f};

    // 获取本地变换矩阵（TR S 复合）
    glm::mat4 GetTransform() const
    {
        glm::mat4 rotation = glm::toMat4(glm::quat(Rotation));
        return glm::translate(glm::mat4(1.0f), Translation)
             * rotation * glm::scale(glm::mat4(1.0f), Scale);
    }
};

struct RelationshipComponent
{
    UUID              ParentID = 0;   // 父实体 UUID，0 = 根节点
    std::vector<UUID> Children;       // 子实体 UUID 列表
};
\end{lstlisting}

两条纪律贯穿全部 20 余个组件：必须有默认构造和拷贝构造（序列化与 Undo
快照依赖值语义）；跨实体引用一律存 UUID 而非指针或 entt::entity
（RelationshipComponent 的 ParentID/Children 就是例证）。UUID 是唯一能
在``保存 $\to$ 重载 $\to$ 快照恢复''全流程中稳定存活的标识符。

\subsection{façade 化：Scene 是门面，不是仓库}

最初的 Scene 是个胖类：注册表、层级逻辑、世界变换计算、运行时资源、
音频视频容器全都塞在里面。三轮重构后它瘦成了六成员 façade
（Engine/src/Scene/Scene.h）：

\begin{lstlisting}
class Scene
{
public:
    Entity CreateEntity(const std::string& name = "Entity");
    void SetParent(Entity child, Entity parent);   // 委托 HierarchyService
    glm::mat4 GetWorldTransform(Entity entity);    // 委托 WorldTransformService
    void OnRuntimeStart();                         // 委托 RuntimeCoordinator
    entt::registry& GetRegistry() { return m_Registry; }
private:
    entt::registry m_Registry;
    SceneEntityIndex m_EntityIndex;          // UUID -> entity O(1) 索引
    WorldTransformCache m_TransformCache;
    ResourceLifecycleCoordinator m_LifecycleCoordinator;
    SceneEnvironmentState m_EnvironmentState;
    std::unique_ptr<SceneRuntimeCoordinator> m_RuntimeCoordinator;
};
\end{lstlisting}

拆出去的逻辑去哪了？变成无状态静态服务类。\texttt{WorldTransformService}
沿 RelationshipComponent 递归上溯父链复合世界矩阵：

\begin{lstlisting}
// WorldTransformService.cpp —— 递归 + 缓存 + 深度防御
glm::mat4 ComputeWorldTransformImpl(reg, entity, index,
                                    cache, depth)
{
    if (depth >= kMaxDepth)   // kMaxDepth=64
    {
        ENGINE_CORE_WARN("... possible cycle in hierarchy");
        return reg.get<TransformComponent>(entity).GetTransform();
    }
    if (cache && cache->TryGet(entity, cached))
        return cached;

    auto& result = reg.get<TransformComponent>(entity).GetTransform();
    if (reg.all_of<RelationshipComponent>(entity))
    {
        auto& rel = reg.get<RelationshipComponent>(entity);
        if (rel.ParentID != 0) {
            entt::entity parent = index.Find(rel.ParentID);
            if (parent != entt::null && reg.valid(parent))
                result = ComputeWorldTransformImpl(
                    reg, parent, index, cache, depth + 1) * result;
        }
    }
    if (cache) cache->Put(entity, result);
    return result;
}
\end{lstlisting}

三个防御细节值得注意：\textbf{kMaxDepth=64 兜底循环层级}（数据损坏时不至
栈溢出）；\textbf{缓存命中短路递归}（编辑器每帧对全场景求世界变换）；而
\texttt{SetParent} 时会先记下子实体的旧世界矩阵，换父后用
$\mathrm{local} = P_{new}^{-1} \cdot W_{old}$ 反解本地变换——保证``拖拽换父
时物体在世界空间里纹丝不动''（Engine/src/Scene/SceneHierarchyService.cpp，
配合 GLM \texttt{matrix\_decompose} 分解回 T/R/S）。

\subsection{双场景快照：Play 按钮的安全网}

点击 Play 会进入物理与脚本的运行态，但编辑器必须能在 Stop 后回到进入前的
精确状态。实现方式不是记录 diff，而是\textbf{整场复制}：EditorSceneSession
维护编辑态与运行态两个 Scene 实例（重构提交 \texttt{b8bceb1}），Play 时把
编辑场景 Copy 一份供模拟使用，Stop 时直接丢弃运行副本。OpenAL Source、
FFmpeg 解码器这类无法安全复制的原生资源由 RuntimeStore 在 Stop 时强制回收
（Phase 1 提交 \texttt{4dab189} 的``运行时资源剥离''正是为此铺路）。
粗暴、简单、可靠——快照成本对一个编辑器会话完全可以接受。

\subsection{Systems 目录：渲染逻辑的静态归拢}

MeshRenderSystem / ShadowSystem / SkyboxSystem / LightSystem /
TerrainRenderSystem / GrassRenderSystem / AudioSystem / VideoSystem 八个
静态系统类构成 Systems/ 目录。它们没有状态、不互相调用，入口全是静态方法，
通过 registry view 批量处理组件。``系统无状态''让渲染管线的执行顺序完全
由 SceneRenderer 编排层决定（第 3 章 8-Pass），职责切分干净。

\section{实现走读:删除一个实体有多难}

DestroyEntity 要处理：递归删除整个子树、清理 UUID 索引、触发资源回调、
从父节点的 Children 列表移除自己。最初实现里``找子树''用的是对每个节点
做全场景扫描，复杂度 $O(N^2)$；优化版（\texttt{641f8ae}）先一次性收集
子树再批量删除。另一个隐蔽 bug 是 UUID 索引允许重复插入
（\texttt{595e461} 修复为检测并拒绝），否则 Find 返回歧义结果，拾取与
序列化都会错乱。EnTT 自身也有兼容性暗礁——\texttt{size\_hint} API 在
版本间行为变化导致实体计数错误（\texttt{2490c50} 改用兼容写法）。

\section{踩坑记录}

\begin{description}
  \item[组件里藏指针] 早期曾有组件直接持有裸指针，Play/Stop 快照恢复后
        变悬垂。教训固化为纪律：组件只存 POD 与 UUID，重资源走 AssetHandle
        （第 6 章）。
  \item[层级环] 手滑把自己设为自己的祖先，世界变换递归立即爆栈。
        kMaxDepth 防御 + IsAncestorOf 前置检查双重保险。
  \item[O(N²) 删除] 大场景删除根节点卡顿数秒，profiler 定位到逐节点
        全表扫描；先收集后删降到 $O(N)$（641f8ae）。
\end{description}

\section{小结与延伸思考}

本章的演化路径——胖 Scene $\to$ façade + 无状态服务 + 协调器——本质是
把``数据（registry）/ 规则（service）/ 流程（coordinator）''三种变化
频率不同的东西分开放。它带来的副产品是可测试性：HierarchyService、
WorldTransformService、EntityIndex 都因此获得了独立单元测试
（tests/ 下 TestSceneHierarchy/TestWorldTransform/TestSceneEntityIndex）。

延伸思考：EnTT 的 sparse-set 对``按类型批量遍历''最优，但对``按实体随机
访问多组件''并不友好。如果未来要做大规模开放世界流式加载，可能需要在
其上再建一层空间索引。当前项目规模下，简单性红利仍然大于替代方案的收益。


\chapter{编译期反射：一套宏驱动的编辑器数据流}
\label{ch:reflection}
% ============================================================
% 第 5 章 编译期反射：一套宏驱动的编辑器数据流
% 锚点：Engine/src/Reflection/ComponentMeta.h, ComponentRegistry.cpp,
%       AutoInspector.cpp, AutoSerializer.cpp, PropertyInfo.h
% 时间线：0399e21(2026-02-25 反射首现) → 83dcf39(Phase 3 元数据统一)
% ============================================================
\section{问题引入：每加一个组件要改五处}

没有反射的日子里，新增一个组件的代价是这样的：在 Components.h 定义结构体
只是第一步；属性面板要手写一段 ImGui 控件；序列化器要手写一对 YAML 读写；
Add Component 菜单要加一项；Undo 系统还要为它写快照逻辑。五处改动、四种
风格、无数漏改的机会——``加了字段但忘了存盘''是当时的高频 bug。

2026-02-25，提交 \texttt{0399e21} 引入轻量级反射系统，此后 Phase 3 重构
（\texttt{83dcf39}）统一组件元数据。最终形态下，新组件只需声明一次元数据，
Inspector 面板、YAML 序列化、Undo 快照三处能力自动到位。

\section{通用原理：C++ 为什么没有反射（以及替代品们)}

截至 C++20，标准 C++ 仍无法在运行期回答``这个结构体有哪些成员''。
工程界的补位方案大致四类：

\begin{itemize}
  \item \textbf{代码生成}：Unreal 的 UHT 在编译前解析宏生成 .generated.h，
        元数据与二进制同源、永远同步，代价是构建管线复杂度陡增
        （参见 docs/ue-reflection-research.md 调研）。
  \item \textbf{运行时注册 + 宏}：宏展开出自描述函数，程序启动时注册进
        全局表。实现轻巧，但成员列表与结构体定义物理分离。
  \item \textbf{第三方库}：rttr、Boost.PFR 等。PFR 需 C++17 结构化绑定
        黑魔法且对嵌套类型支持有限；rttr 引入重依赖。
  \item \textbf{等标准}：C++26 静态反射提案仍在路上。
\end{itemize}

本项目选了第二类——对一个毕业设计体量的项目，``不引入构建步骤、不引入
依赖、半小时能讲清楚''的权重远高于``成员列表零分离''的美学诉求。

\section{本项目的设计与决策}

\subsection{宏的三层分工}

反射宏分三层，各司其职（Engine/src/Reflection/ComponentMeta.h）。
第一层声明组件身份：

\begin{lstlisting}
// 声明组件显示名 —— 展开为命名空间 + inline 函数
#define ENGINE_COMPONENT(CompType, displayName)              \
    namespace _Reflect_##CompType                            \
    {                                                        \
        inline ::Engine::ComponentMeta& GetMeta()            \
        {                                                    \
            static ::Engine::ComponentMeta s_Meta;           \
            s_Meta.TypeName    = #CompType;                  \
            s_Meta.DisplayName = displayName;                \
            s_Meta.Properties.clear();                       \
            return s_Meta;                                   \
        }                                                    \
    }
\end{lstlisting}

关键技巧：\texttt{\#CompType} 把模板参数字符串化，每个组件获得独立的
\texttt{\_Reflect\_XXX} 命名空间——同名属性互不冲突。\texttt{inline}
函数内的 static 局部变量保证全进程单例。

第二层逐个登记属性：

\begin{lstlisting}
#define ENGINE_PROPERTY(CompType, field, displayName, propType) \
    namespace _Reflect_##CompType {                             \
      namespace _Prop_##field {                                 \
        inline void Register(::Engine::ComponentMeta& meta)     \
        {                                                       \
            ::Engine::PropertyInfo info;                        \
            info.Name        = #field;                          \
            info.Type        = PropertyType::propType;          \
            info.Offset      = offsetof(CompType, field);       \
            meta.Properties.push_back(info);                    \
        }                                                       \
      }                                                         \
    }
\end{lstlisting}

\texttt{offsetof} 是整套方案的核心魔法：运行期拿到成员偏移量后，
\texttt{(char*)component + offset} 即可无类型地访问任意组件的任意字段——
这正是 AutoSerializer 与 AutoInspector 都不需要知道具体类型的底气。

第三层把一切组装并触发注册（.cpp 中使用）：

\begin{lstlisting}
// ComponentRegistry.cpp 中 —— 以 RigidBodyComponent 为例
REGISTER_COMPONENT_BEGIN(RigidBodyComponent)
    REGISTER_COMPONENT_PROPERTY(RigidBodyComponent, BodyType)
    REGISTER_COMPONENT_PROPERTY(RigidBodyComponent, Mass)
    REGISTER_COMPONENT_PROPERTY(RigidBodyComponent, LinearVelocity)
REGISTER_COMPONENT_END(RigidBodyComponent)

// END 展开的关键尾部：
//   ::Engine::ComponentRegistry::Instance().Register(meta);
//   return true; } ();
\end{lstlisting}

\texttt{BEGIN} 展开成 \texttt{static bool s\_Registered\_X = []() -> bool \{}
——一个立即执行的静态初始化 lambda。它在 main() 之前运行，顺手把七个
类型擦除操作挂上 meta：Has/Add/Get/Remove/Copy 五个操作 EnTT registry，
而 \textbf{Snapshot/Restore} 用 \texttt{std::any} 携带整个组件 POD 的值拷贝
——这两个槽位正是第 12 章 Undo 系统的钩子。反射系统在这里埋下了
编辑器撤销能力的地基。

\subsection{一次遍历，两处消费}

消费端是两个高度对称的遍历器。AutoSerializer 按 PropertyType 分支写 YAML：

\begin{lstlisting}
// AutoSerializer.cpp —— 反射驱动序列化
void AutoSerializer::Serialize(YAML::Emitter& out,
                               const ComponentMeta& meta, void* component)
{
    for (auto& prop : meta.Properties)
    {
        if (prop.Hints.Transient) continue;   // 跳过瞬态字段
        void* ptr = static_cast<char*>(component) + prop.Offset;
        switch (prop.Type) {
        case PropertyType::Float:
            out << YAML::Key << prop.Name << YAML::Value
                << *static_cast<float*>(ptr);
            break;
        // Int/Bool/String/Vec2-4/Enum ... 同构分支
        }
    }
}
\end{lstlisting}

AutoInspector 几乎是它的镜像——同样的循环、同样的 switch，只是 case 里
换成 ImGui::DragFloat / Checkbox / DragFloat3。这种对称不是巧合：
\emph{属性面板和存档文件本来就是同一份元数据的两种投影}。hints 系统
（Min/Max/Speed/Format/Group/Transient，PropertyTypes.h）进一步让 UI 行为
（拖动速度、范围钳制）与序列化行为（Transient 跳过）都由声明处一处控制。

\subsection{策略标志：自动化留后门}

并非所有组件都适合全自动。\texttt{CustomUI}/\texttt{CustomSerial} 两个
Flag 允许组件保留手写 Inspector 或手写序列化（如 MeshRendererComponent
的材质预览需要特殊绘制）。自动化的价值在于覆盖那 80\% 规整组件，
而不是强行吞掉剩下 20\% 的特例。

\section{实现走读:注册表的运行时形态}

ComponentRegistry 是个单例 map：TypeName $\to$ ComponentMeta。编辑器的
Add Component 菜单遍历它列出所有可加组件；属性面板按实体的组件集合查表
绘制；场景反序列化时用 TypeName 字符串找回 Meta 并驱动 Deserialize。
一条链路串起三个子系统——这也是``未注册的组件不会出现在编辑器中也不会被
序列化''这条规则（.claude/rules/reflection.md）的原因：没进表就没有
任何消费者认识你。

\section{踩坑记录}

\begin{description}
  \item[PropertyType 与实际类型不匹配] 宏只记录声明的枚举，若某字段实际
        是 double 却声明成 Float，offsetof 照常工作但读出的字节被错误
        解释——画面错乱且难排查。纪律：声明即契约，改字段必须同步改宏。
  \item[Enum 属性忘配 EnumNames] Enum 类型需要外部提供名字数组并设置
        hints.EnumCount，漏掉后面板下拉框为空。
  \item[u8 字面量陷阱] C++20 下 \texttt{u8"中文"} 是 char8\_t 数组，
        无法传给 ImGui 的 const char* 参数；项目以 MSVC /utf-8 +
        普通字面量绕开（迁移文档 P-7 同款教训）。
\end{description}

\section{小结与延伸思考}

这套约 150 行宏代码换来的杠杆率惊人：20 余个组件、数百个属性的
UI+序列化+快照全部声明式搞定，新增组件的心智负担从``改五处''降到
``写三行宏''。它是编辑器（第 11 章）、Undo（第 12 章）两个后续章节
共同的基石。

延伸思考：宏方案的真正成本在``声明与定义分离''——结构体改了字段名，
宏这边的字符串不会报错，只能靠 code review 拦截。C++26 静态反射落地后，
同样的架构可以用 constexpr 元编程重写，届时 offsetof 与字符串化都将
成为历史。


\chapter{资产管理：SlotMap、异步加载与热重载}
\label{ch:asset}
% 第 6 章 资产管理（占位）
锚点：\texttt{Asset/SlotMap.h}（世代槽资源池）、\texttt{AssetHandle.h}
（强类型句柄）、\texttt{AssetManager.cpp}、\texttt{AsyncLoadQueue}
（后台加载）、\texttt{FileWatcher}（热重载）。

素材线索：提交 \texttt{3955e2e}``AssetManager 完整实现，SlotMap 资源池 +
异步加载 + 热重载''；SlotMap::Remove 屏蔽保留槽防 freelist 污染
（\texttt{d05ea6b}）；AssetHandle 加 Type 字段根治 SlotMap 冲撞断言
（\texttt{182f5e7}）；线程安全 LoadAsync 去重（\texttt{576116a}/\texttt{6ef37c9}）。

% TODO(M3): SlotMap 世代计数原理；句柄 vs 裸指针；异步队列去重。


% ---- 第 III 部分：让世界可信 ----
\chapter{物理系统：从手写冲量到 Bullet3}
\label{ch:physics}
% ============================================================
% 第 7 章 物理系统：从手写冲量到 Bullet3
% 锚点：Engine/src/Physics/BulletPhysicsWorld.h/.cpp, SDFMath.h
% 时间线：7341800(Phase 9 手写+Bullet) → 89fe5b4(审计10bug)
%         → 9cff9ce(移除自研强制Bullet)
% ============================================================
\section{问题引入：两个方块穿过了彼此}

物理是``不做不知道难，做了才知道有多难''的模块。第一版手写物理上线当天
就能复现经典惨案：快速移动的方块一帧位移超过对方厚度，穿透检测失效；
斜面上的盒子抖动不止；堆叠的箱子越叠越陷。审计报告一次性列出十项 bug
（修复提交 \texttt{89fe5b4}），其中``手写物理穿模''位列榜首。

本章讲一个诚实的故事：项目先手写了一遍冲量法物理（Phase 9，
\texttt{7341800}），在理解了它难在哪里之后，把动力学部分整体换成了
Bullet3（\texttt{9cff9ce} ``移除 Custom PhysicsWorld，强制使用 Bullet3''）——
但手写阶段的碰撞数学遗产至今仍在服役。

\section{通用原理：冲量法动力学}

刚体物理的主流解法是顺序冲量法（sequential impulse）：每个约束（接触、
关节）迭代求解冲量 $\lambda$，使相对速度满足约束方程；位置误差用 Baumgarte
稳定项修正。要让它工作，至少需要：宽相检测（AABB 树/扫掠剪枝）、窄相检测
（精确几何相交与接触点生成）、接触流形管理（防止抖动）、摩擦与恢复系数
模型、惯性张量（旋转响应的正确性全靠它）。

每一项单看都不复杂，合起来就是数万行成熟库代码。Bullet3 的价值正在于此：
它的 \texttt{btSequentialImpulseConstraintSolver} 是这套方法二十年的工程
结晶。

\section{本项目的设计与决策}

\subsection{双轨制：ECS 与 Bullet 世界并行}

BulletPhysicsWorld 不接管实体，而是作为 ECS 的``影子世界''运行
（Engine/src/Physics/BulletPhysicsWorld.h）：

\begin{lstlisting}
// 碰撞事件 —— 面向 gameplay 的三态语义
enum class CollisionEventType { Enter, Stay, Exit };

struct CollisionEvent
{
    CollisionEventType Type;
    entt::entity EntityA, EntityB;
    glm::vec3 ContactPoint;
    glm::vec3 ContactNormal;   // A -> B 方向
    float     Impulse;
    bool      IsTrigger;
};

class BulletPhysicsWorld
{
public:
    void Init(glm::vec3 gravity = {0, -9.81f, 0});
    void Step(float dt, entt::registry& reg,
              const SceneEntityIndex& index);
    void CreateBodies(entt::registry& reg,
                      const SceneEntityIndex& index);
    void SyncToECS(entt::registry& reg,
                   const SceneEntityIndex& index);   // 模拟结果回写
    void SyncFromECS(entt::registry& reg,
                     const SceneEntityIndex& index); // Kinematic 同步
};
\end{lstlisting}

每帧流程：CreateBodies 补建新实体的 btRigidBody → Step 推进仿真 →
SyncToECS 把 Bullet 的变换写回 TransformComponent。事件系统维护上一帧
接触对集合，差分出 Enter/Stay/Exit 三态——这是 gameplay 脚本最想要的
接口形态，而不是 Bullet 原生的回调风暴。

\subsection{碰撞体包装：偏移量与复合形状}

组件里声明的 ColliderOffset（碰撞体相对实体原点的偏移）不能直接喂给
Bullet——btRigidBody 的质心变换有自己的约定。项目用 \texttt{btCompoundShape}
包装子形状并设置子形状局部偏移来承载这一信息（修复提交 \texttt{a4aa1ad}，
第 4 章的 UUID 索引在这里再次发挥作用：ECS 实体到 Bullet body 的映射靠
UUID 稳定锚定）。地形则用 \texttt{btTriangleMesh} 从高度图网格构建，
BodyInfo 聚合初始化还曾因漏掉 childShape 字段编译失败
（\texttt{adbe4a8} 小修）。

\subsection{手写遗产：仍然服役的碰撞数学}

移除 Custom PhysicsWorld 删掉的是动力学求解器；碰撞几何数学被完整保留。
Box-Box 从 AABB 近似升级为 OBB-OBB SAT（分离轴定理，
\texttt{1e3a12e}）、接触点从两中心中点改为真实接触面计算
（\texttt{31fefb5}）、球-球零距离边界条件（\texttt{bff156f}）——这些
修复沉淀为 SDFMath.h 与碰撞数学工具。SAT 的退化轴处理、惯性张量变换公式
甚至有专门的数学文档（docs/ 下 CollisionMath 与 PhysicsWorld 笔记）。

SDFMath.h 展示了``可测性提取''的典型手法：

\begin{lstlisting}
// 纯数学函数 —— header-only，兼容 CUDA 和 C++。
// 从 CudaSPHPipeline.cu 提取以便单元测试。
#ifdef __CUDACC__
#define SDF_DEVICE __device__ __host__
#else
#define SDF_DEVICE
#endif

namespace Engine { namespace SDFMath {
    struct SDFResult { float sdf; float normalX, normalY, normalZ; };

    // Box SDF: 点到 AABB 的有符号距离（局部坐标系）
    SDF_DEVICE inline SDFResult BoxSDF(float lx, float ly, float lz,
                                       float hex, float hey, float hez);
}}
\end{lstlisting}

同一份头文件同时服务 C++ 单元测试与 CUDA kernel（第 9 章 SPH 流体的
网格碰撞用它），一份数学两处消费，测试只需编译 C++ 侧。

\section{实现走读:旋转的三个深坑}

手写阶段留下的最有价值的教训集中在旋转积分上，每一条都有对应修复提交：

\begin{itemize}
  \item \textbf{万向节锁}：用欧拉角积分角速度，特定姿态下自由度丢失。
        修复（\texttt{e45b642}）：改持久四元数积分——四元数无奇点，
        且必须持久保存而非每帧由欧拉角重建。
  \item \textbf{陀螺力矩}：角动量守恒要求 $\dot{L} = \omega \times L$，
        漏掉这项会让非对称刚体旋转漂移（\texttt{ecf3017} 补上陀螺力矩项）。
  \item \textbf{惯性张量}：标量惯性矩对旋转刚体根本不够，升级为 mat3 并
        加平行轴定理修正（\texttt{b85497d}，配套旋转/缩放测试覆盖）。
\end{itemize}

这三个坑共同指向同一句话：\emph{3D 旋转没有免费午餐}。2D 物理里标量角速度
+标量惯量的直觉在 3D 全部失效。

\section{踩坑记录}

\begin{description}
  \item[双重位置矫正] 手写阶段对同一接触先后施加两次位置修正，堆叠场景
        能量凭空增加导致``爆炸''。\texttt{e45b642} 一并移除。
  \item[外力清零时机] 施加的外力若不在 Step 末尾清零，会持续累积下一帧
        加速度——表现为物体``自己加速''。
  \item[负缩放] 变换矩阵含负缩放（镜像）时 AABB 尺寸计算出负值，
        宽相直接失明，需 abs 处理（\texttt{7dc27ef} 系列 Phase 1 修复）。
  \item[世界坐标 vs 本地坐标] 碰撞检测曾误用本地坐标比较不同父级下的
        实体（\texttt{ad0bf7d} 改世界坐标）——与第 8 章 FluidPass 的
        世界坐标 bug 同源，见该章踩坑录。
\end{description}

\section{小结与延伸思考}

物理系统的叙事弧线——手写、碰壁、换库、留其精华——本身就是一堂工程课：
\emph{造轮子的目的不是用上这个轮子，而是理解轮子}。最终架构里，Bullet 负责
动力学求解（信任其二十年打磨），自研代码负责 ECS 桥接、事件语义、以及
SPH 流体需要的 SDF 碰撞查询——后者恰是 Bullet 不擅长而本项目刚需的部分。

延伸思考：如果重做，会一开始就用 Bullet 吗？大概率不会。正是手写阶段对
穿透、抖动、万向锁的亲历，让后来读 Bullet 文档时每个参数都``认识''——
restitution 怎么调、ERP/CFM 在解决什么，不再是抽象名词。教学价值与工程
效率的取舍，毕业设计显然应该偏向前者。


\chapter{SPH 流体与 GPU 粒子：Compute 管线三部曲}
\label{ch:sph-gpu}
% ============================================================
% 第 8 章 SPH 流体与 GPU 粒子：Compute 管线三部曲
% 锚点：ParticleSystemGPU, FluidSystemGPU, SpatialHashGrid,
%       SPHKernelMath.h, GPUAsyncReadback + 14 个 sph/fluid/grid shader
% 时间线：29d41cc/beb5ac0(Phase 10) → 58a69bf(PCISPH)
%         → e46f088(SSFR) → 6cf61fb(Akinci 表面张力)
% ============================================================
\section{问题引入:一万颗粒子的 CPU 瓶颈}

粒子系统初版跑在 CPU 上，五万粒子时单帧模拟耗时超过 30ms——还没渲染就
已经跌破 30 FPS。而 GPU 正闲着。Phase 10 系列（\texttt{29d41cc} 基础设施
与 4-Pass 管线、\texttt{beb5ac0} 空间哈希网格与 SPH 流体）把整个模拟搬上
GPU Compute Shader，随后一路演进：PCISPH（\texttt{58a69bf}）、Screen-Space
Fluid Rendering（\texttt{e46f088}）、Akinci 表面张力（\texttt{6cf61fb}）。
本章是全书技术密度最高的一章，也是踩坑最多的一章。

\section{通用原理:SPH 与邻域搜索}

SPH（Smoothed Particle Hydrodynamics）把流体离散为携带质量/速度/密度的
粒子，任意场量用光滑核 $W(r,h)$ 加权求和近似：
$\rho_i = \sum_j m_j W(\lvert \mathbf{r}_i - \mathbf{r}_j\rvert, h)$，
压力由状态方程给出，加速度含压力梯度、粘性与外力。WCSPH 每步直接解压力；
PCISPH（Solenthaler 2009）则迭代预测位置修正密度误差，允许更大时间步。
本项目两者都实现，共享同一套核函数数学（SPHKernelMath.h，header-only 可测）。

朴素邻域查找是 $O(N^2)$——对每对粒子算距离。空间哈希把它降到近似
$O(N)$：粒子按格子哈希分桶，只查邻居桶。GPU 实现的经典三段式是
hash $\to$ prefix sum $\to$ scatter。

\section{本项目的设计与决策}

\subsection{Compute 4-Pass 粒子管线}

ParticleSystemGPU 把粒子生命周期拆成四个 compute dispatch：

\begin{itemize}
  \item \textbf{emit}：按发射参数生成新粒子写入粒子池；
  \item \textbf{simulate}：积分速度与位置（可挂接 SPH 力或 Bullet 刚体碰撞）；
  \item \texttt{render\_args}：在 GPU 上统计存活数并填充 Indirect Draw 参数块；
  \item \textbf{billboard}：顶点着色器按粒子属性展开面向相机的四边形。
\end{itemize}

关键决策是第 3 步：存活计数与绘制参数完全在 GPU 上生成，CPU 全程不知道
具体粒子数，DrawArraysIndirect 一条命令完成提交——CPU-GPU 数据往返被
压缩到零。VMware SVGA3D 不支持 indirect 时回退实例化绘制（renderer.md
规则记录的兼容性分支）。

\subsection{SpatialHashGrid:三段式与 binding 布局}

Engine/src/Renderer/SpatialHashGrid.h 是网格的完整自述：

\begin{lstlisting}
// Grid SSBOs —— binding 编号是跨 shader 的全局契约
Ref<ShaderStorageBuffer> m_CellHash;      // per-particle hash (binding 1)
Ref<ShaderStorageBuffer> m_CellCount;     // cell histogram   (binding 6)
Ref<ShaderStorageBuffer> m_CellStart;     // prefix sum 结果  (binding 5)
Ref<ShaderStorageBuffer> m_SortedIndices; // 重排索引         (binding 7)

void Build(uint32_t aliveCount, bool usePredictedPos = false);
// Vulkan 主帧 cmd 录制版本（避免 SingleTime stall，见 SPEC D-3）
void BuildVulkan(VkCommandBuffer cmd, uint32_t aliveCount,
                 bool usePredictedPos = false);
\end{lstlisting}

prefix sum 本身又是三个 pass（块内扫描 / 块和扫描 / 块偏移加回），pass 间
必须插 memory barrier——省略任何一个都会读到脏数据。binding 编号表是
14 个 sph/fluid/grid shader 共同遵守的全局布局契约，它的脆弱性本章踩坑录
有血泪案例。

\subsection{PCISPH 与表面张力:站在文献肩膀上}

PCISPH 的 $\delta$ 参数没有手抄论文数值，而是按 Solenthaler 2009 的公式
在初始化时自动推导（\texttt{00146b7}）；表面张力实现 Akinci 2013 的曲率
修正项，且 GLSL 与 CUDA 双实现同步落地（\texttt{6cf61fb}）。学术参数的
``自动推导优于魔法数字''在这里得到验证：改光滑核半径时一切自适应。

\subsection{SSFR:从粒子到水面}

粒子本身不像流体。SSFR（Screen-Space Fluid Rendering）管线用五个额外
pass 把粒子变成液体：thickness（厚度累积染色）$\to$ depth（椭圆 splat
写最小深度）$\to$ bilateral smooth（深度域平滑）$\to$ 从深度重建法线
$\to$ composite（折射+反射+泡沫合成进主场景）。它牺牲了物体空间一致性
换取巨大性能优势——屏幕空间的平滑天然只在可见处花钱。

\section{实现走读:异步回读与 fence}

CPU 需要知道存活数（如发射器停止条件），最初实现每帧
\texttt{glGetBufferSubData} 同步回读——GPU 必须等所有在途工作完成才能
服务这个读取，实测造成 28ms 的每帧 stall（修复提交 \texttt{01de4e3}）。
方案是 GPUAsyncReadback 抽象：N 个 host-visible 缓冲轮转，第 $k$ 帧发起
copy、读第 $k-N$ 帧的结果，fence 保证就绪。代价是数据延迟 N 帧——对
``上一帧存活数''这种用途完全无感。这次教训后来直接塑造了 Vulkan 路径的
D-11 决策（3 槽 ring + 独立 fence，见第 14 章）。

\section{踩坑记录}

这一章的坑值得单独成节——每个都是 GPU 并发世界的经典陷阱：

\begin{description}
  \item[SSBO slot 被 Grid.Build 覆盖] PCISPH 迭代中途调用 Grid.Build()
        重建网格，后者绑定的 SSBO 恰好占了 slot 1——预测位置缓冲被顶掉，
        流体整场不可见且不报任何错（\texttt{e63d819}）。排查花了三天，
        最终靠逐 pass 二分绑定才定位。\emph{OpenGL 全局 binding 空间里，
        ``谁绑了什么''是隐式契约}。
  \item[GL\_BLEND 泄漏] SSFR 深度 pass 继承了上游未关闭的混合状态，
        最小深度混合成了``平均深度''，水面法线重建全错（\texttt{b85f9d5}
        在深度与平滑通道显式禁 blend）。全局状态机的又一记闷棍。
  \item[CUDA/GLSL 一致性] 密度自贡献项遗漏、预测位置恢复时机、积分顺序
        三处差异让两后端结果分道扬镳（\texttt{dc72895} 三连修）。浮点
        求和顺序不同本属正常，但算法级差异必须消灭——这正是附录 A 容差
        协议存在的原因。
  \item[差一帧的 alive list] emit 之后不重建 alive list 就 simulate，
        新粒子晚一帧参与模拟（\texttt{547e3ad}）。GPU 数据结构没有
        ``自动维护的不变量''，每次写入都要想清楚下游谁在读。
  \item[cellSize=2h 的冗余] 早期邻居搜索半径设 2h，27 邻域扫描近半无效；
        收紧为 h 后消除冗余（\texttt{7c6c746}）。理论最优与实现常量
        要对齐验证。
\end{description}

\section{小结与延伸思考}

三部曲的主线：4-Pass 管线确立 CPU 极简参与；三段式网格把邻域查询摊平到
GPU；文献算法（PCISPH/Akinci）以``自动推导 + 双端一致''方式落地；SSFR
补上最后一公里的观感。性能上，这套管线支撑了附录 A 五万粒子的正式基准。

延伸思考：屏幕空间流体在粒子被遮挡或飞出屏幕时会丢失厚度信息（边缘
``融化''伪影）；物体空间方法（如 Marching Cubes 提取 mesh）能根治但代价
高昂。另一个值得探索的方向是把 PCISPH 迭代也做进 indirect 结构，进一步
减少 CPU 参与——目前迭代次数仍是启动时常量。


\chapter{CUDA sidecar：引入、移除与重引入}
\label{ch:cuda-sidecar}
% 第 9 章 CUDA sidecar（占位）
锚点：\texttt{Platform/CUDA/} 11 文件约 2370 行——
\texttt{CudaGLInteropContext}、\texttt{CudaParticlePipeline.cu}/
\texttt{CudaSPHPipeline.cu}、\texttt{CudaTimingHelper}、
\texttt{CudaErrorHandling.h}、\texttt{CudaPoisonState.h}
（进程级 atomic 中毒标志 + 永久回退 OpenGL Compute 路径）。

三幕剧叙事（构建历程最佳素材）：引入（\texttt{7290534} 起 CUDA 四部曲
调研文档）→ 移除（\texttt{e71ea47}/\texttt{a18fd0c}``移除全部 CUDA 代码''
23 文件 4566 行）→ 重引入（\texttt{a1f9def} 起 wip 系列 + Pimpl 隐藏
CUDA 依赖 \texttt{6969f20}/\texttt{708eba8} + 数据布局静态断言
\texttt{1fd159d}）。

% TODO(M3): 互操作内存共享机制；中毒回退状态机（有单元测试）；A/B
% 自动交替基准。


\chapter{音视频与地形脚本：外围模块速写}
\label{ch:peripherals}
% 第 10 章 外围模块速写（占位）
音频（OpenAL，590 行）、视频（FFmpeg RTSP/文件解码，570 行）、地形
（高度图 + SplatMap PBR，300 行）、脚本（NativeScript + Lua 双后端，
1708 行，含热重载）、事件系统、性能监控。

素材线索：Phase A Audio + RTMP Streaming（\texttt{8eca5f6}）；Lua 后端
元数据接入（\texttt{b36119b} 起 6 连击）；ScriptEditorPanel 与代码补全
方案 B 设计文档。

% TODO(M3): 每模块一节短文，重点讲"引擎如何为它们让位"（生命周期协调）。


% ---- 第 IV 部分：编辑器与工程化 ----
\chapter{编辑器架构：EditorLayer 的膨胀与瘦身}
\label{ch:editor}
% 第 11 章 编辑器架构（占位）
锚点：\texttt{EditorLayer.h/.cpp}（最大文件，总协调器）、
\texttt{EditorSceneSession}（场景会话 + Play/Stop 双场景模型）、
\texttt{EditorShell}（菜单快捷键）、\texttt{EditorPanelCoordinator}（七面板）、
\texttt{EditorRenderController}、\texttt{EditorViewportController}（拾取映射）、
\texttt{EditorSelectionGizmoController}（ImGuizmo）。

叙事线：2026-03 连环重构——十余个 refactor commit 将上帝类 EditorLayer
逐块拆出（提取 RenderSettingsPanel → SceneSession → Shell → Viewport →
Gizmo → PanelCoordinator），docs/plans/refactor-* 五份计划 + 收口总结。

% TODO(M3): 上帝类是怎么形成的；拆分顺序与依赖方向；崩溃转储钩子。


\chapter{Undo/Redo：命令模式与反射快照}
\label{ch:undo}
% 第 12 章 Undo/Redo（占位）
锚点：\texttt{Editor/UndoSystem.h/.cpp} + \texttt{CommandHistory.h/.cpp}
（执行/记录语义分离，14 单元测试用例）。

机制：命令模式 + 反射系统 Snapshot/Restore（\texttt{std::any} 携带组件
POD 快照）。踩坑素材：EntityCreateCommand Redo 保留原始 UUID
（\texttt{fd1eaee}）；EntityDeleteCommand 递归保存/恢复子树
（\texttt{27f6c4c} 系）；多选变换 undo（\texttt{62479f9}）；Play 模式
暂停命令记录（\texttt{2371ecf}）。

% TODO(M3): 命令 vs 快照两种 Undo 流派对比；与反射系统的耦合设计。


\chapter{性能工程：三次真实调优案例}
\label{ch:performance}
% 第 13 章 性能工程（占位）
三个完整案例（docs/perf/ 三部曲，均有方法论 + 前后数据）：
(1) VTune Hotspots → uarch 定位鼠标卡顿：GetFileAttributesExW 文件系统
热点，修复后对比（mouse-lag 文档）；
(2) NSight + CSV 交叉分析 50k 粒子 SPH 帧时间，VSync 量化阶梯根因
（round3 文档）；
(3) uarch 探索：Retiring 仅 13.8%、Back-End 47% 瓶颈 + CUDA 同步自旋
分析（vtune-optimization-guide）。

代码锚点：\texttt{Debug/PerformanceMonitor}（DrawCalls/Triangles 统计
内联在 RenderCommand）、\texttt{ProfileTimer}、\texttt{GPUTimerQuery}
（GL + CUDA 双路径计时埋点 \texttt{aa38854}）。

其他素材：spin-wait 限帧 56%→0%（\texttt{9eb9a79}）；immutable storage
+ SPH kernel 预计算 + Bloom 半分辨率（\texttt{d9d9168}）；
AssetBrowserPanel 每帧遍历文件系统修复（\texttt{ce1e31c}）；
normalMatrix GPU→CPU（\texttt{a328224}）。

% TODO(M3): 每案例一节：症状→工具→定位→修复→数据。


\chapter{Vulkan 后端迁移：一次受控的双后端改造}
\label{ch:vulkan}
% 第 14 章 Vulkan 后端迁移（占位——全书最重章，素材最富）
锚点：\texttt{Platform/Vulkan/} 40 文件约 6000 行（Context/Swapchain/
Allocator(VMA)/RenderPass/Pipeline/PipelineCache/Descriptor 三层/Shader
(shaderc+spirv-cross)/Buffer/Texture/VertexArray/Synchronization/
BarrierUtil/AsyncReadback/IBLGenerator）。

叙事骨架（SPEC.md Phase 1-8 即章节骨架）：
Phase 1 抽象泄漏归零（前置工程）→ Phase 2 基础设施清屏 →
Phase 3 RendererAPI 核心 → Phase 4 shaderc 运行时编译 →
Phase 5 VMA 资源族 → Phase 6 ImGui 多视口 → Phase 7 Compute 全家桶迁移
（IBL/草地/网格/粒子/SPH）→ Phase 8 PBR 主路径接通（进行中，
DrawIndexed 待实装）。

决策素材：SPEC §3 十六条 D-x 决策全部按``背景+结论+影响范围''格式现成；
ADR-0001（shader 单源双路径三方案对比）/ADR-0002（主帧 cmd 暴露策略）
标准 MADR 格式；Pitfalls 表 17 条即本章``踩坑录''小节初稿。

% TODO(M3): 与 OpenGL 抽象层的映射表；双路径 shader 宏策略；验证方法学
%（validation layer + 冒烟层）。


% ---- 附录 ----
\appendix
\chapter{三后端流体基准实验协议}
% ============================================================
% 附录 A 三后端流体基准实验协议
% 素材：docs/fluid-compute-benchmark-spec.md
%       docs/thesis/experiment-data-summary.md
% 提交群：37daaef..3cfee23（约 30 个 benchmark 工程 commit）
% ============================================================
\section{为什么需要一个``基准协议''}

第 8、9、14 章分别讲了 OpenGL Compute、CUDA sidecar 与 Vulkan Compute
三条流体计算路径。它们各自``能跑''不等于``可比''——CUDA 路径曾经自动
交替 A/B 运行（\texttt{0d5c4fd}），拿这种数据写论文等于自欺。要让三后端
的对比站得住脚，需要一份事先冻结的实验协议，而不是跑完数据后再挑好看的。

本附录整理 docs/fluid-compute-benchmark-spec.md 固化的协议要点与工程实现，
它服务于毕业论文第五章，也是本书所有性能结论的数据来源。整个体系由约
30 个提交逐步建成（\texttt{37daaef} 起），每一步解决一个方法论问题。

\section{公平性约束：先定义什么叫``同一场比赛''}

协议的核心条款：

\begin{itemize}
  \item \textbf{确定性输入}：三后端必须使用同一份初始粒子状态，固定随机
        种子（42）、时间步长（1/120s）与求解器参数；禁止依赖墙钟时间生成
        随机状态（\texttt{66b3117} 实现确定性初始化）。
  \item \textbf{等步数比较}：每后端执行固定 120 个不计时验证帧后才能
        比较，杜绝``步数不一致导致密度天然不同''的假差异
        （\texttt{9dd4274} 校验等步数一致性）。
  \item \textbf{独立进程}：每个后端单独进程运行，禁止把逐帧交替 A/B 的
        结果当论文数据。
  \item \textbf{计时边界冻结}：GPU Compute 时间只覆盖网格/密度/力/积分
        四阶段；初始化、正确性回读、CSV 写入不计入。CUDA 的 CUB Prefix
        Sum 必须算进 Compute，而 CUDA-GL Map/Unmap 的主机调用单列为
        Interop 时间。
  \item \textbf{回退即失败}：后端中毒或回退时该组实验判失败，禁止用
        回退后端的数据冒充目标后端。
\end{itemize}

\section{实验矩阵与命令行接口}

固定矩阵三维展开：后端（OpenGL/CUDA/Vulkan）$\times$ 求解器（WCSPH/
PCISPH，PCISPH 固定 6 迭代）$\times$ 粒子规模（1k/5k/10k/20k/50k），
每组预热 100 帧、采样 1000 帧、独立重复 5 次。全部参数经命令行显式传入并做范围校验，无效组合返回非零退出码：

\begin{lstlisting}[language=bash]
Editor.exe --benchmark-fluid --backend opengl --solver pcisph \
  --particles 10000 --iterations 6 --warmup 100 \
  --frames 1000 --runs 5 --fixed-dt 0.008333333 \
  --seed 42 --output benchmark/results.csv
\end{lstlisting}

工程侧配套三件套：断点续跑（\texttt{328d936}，已完成组跳过，矩阵中断后
可恢复）、实验矩阵自动执行与统计汇总（\texttt{738f7e2}）、互操作与端到端
耗时汇总（\texttt{3f1acce}）。原始逐帧 CSV 必须保留——协议明文禁止只存
汇总表，这是可复现性的底线。

\section{数值容差：从校准数据到冻结阈值}

不同着色器编译器、不同浮点执行顺序决定了三后端不可能逐位一致。容差不能
拍脑袋，也不能看结果再定（那是 p-hacking）。协议采用两步法：

\begin{enumerate}
  \item \textbf{短矩阵校准}：2026-07-19 用 1000 粒子跑六组三后端对照，
        相对 OpenGL 基线的平均密度最大相对偏差为 0.0261\%（CUDA/PCISPH）。
  \item \textbf{冻结容差}：正式矩阵前冻结平均密度相对容差 0.1\%
        （约为校准最大偏差的 3.8 倍，\texttt{ea56c76} 固化）。它足以容纳
        编译器差异，又能拒绝此前由压力语义错误造成的百分比级算法偏差。
\end{enumerate}

任一组超容差即标记正确性失败、不计算加速比——宁可缺数据，不可假数据。

容差条款不是凭空设防。三后端曾存在真实的语义分歧：压力求解的积分方式在
GLSL 与 CUDA 实现中不一致，密度结果差出百分比级。排查与对齐由修复提交
\texttt{c801216} 完成（对齐三后端流体求解语义），PCISPH 的位置积分随后也
做了同样对齐（\texttt{94f6091}）。\emph{先修语义、再定容差、后跑矩阵}——
这个顺序反过来就是灾难。

\section{统计规则}

丢弃全部预热帧；保留正常慢帧、不按主观阈值删离群点（只剔除查询未就绪/
设备错误等明确无效样本）；报告均值、标准差、中位数、P95、极值；加速比
定义为基线均值除以目标均值。协议特别强调：论文同时报 GPU Compute 时间与
端到端时间，禁止用不含 Swap/Present 的引擎帧时间冒充真实显示帧时间。

正式数据快照见 docs/thesis/experiment-data-summary.md；可视化图表由
docs/thesis/figures/scripted/ 下的 Python/uv 脚本从原始 CSV 再生成
（generated/ 目录含公式示意与实验结果共 17 组 PDF/PNG），保证``图能从
数据重算出来''。

\section{复现指南}

在本仓库复现全部实验的最短路径：(1) 按各后端要求构建（Vulkan 需 SDK
1.4+ preset，CUDA 需 ENGINE\_ENABLE\_CUDA=ON）；(2) 依次执行上节命令行
形式的 $3\times2\times5=30$ 组参数；(3) 用汇总脚本生成统计表；(4) 对照
experiment-data-summary.md 的冻结参数核验场景配置未漂移。断点续跑机制
让第 (2) 步可以随时中断重入。

\section{延伸思考}

这份协议最大的经验是：\emph{基准的价值在跑之前就决定了}。确定性种子、
计时边界、容差冻结、回退判失败——每一条都是在``数据出来之后想改规则''
的冲动被制度提前锁死。这套纪律适用于任何性能声明场景，与具体引擎无关。


\backmatter
% \bibliography{references}

\end{document}
